% source: J. S. Milne, "Abelian Varieties" (course notes), v2.00, 2008, www.jmilne.org
% pdf_page: 0022
% doc_page: 16 (Chapter I, Section 3 — "Rational maps."; Theorem 3.1 and the start of its proof)
% retrieved: 2026-07-17
% transcribed_by: claude-fable-5 (vision, rendered at 200dpi via pdftoppm)

% [running head: CHAPTER I. ABELIAN VARIETIES: GEOMETRY, page 16]

% (no non-standard operators on this page; rational-map dashed arrows rendered with \dashrightarrow from amssymb)

\subsection*{Rational maps.}

We first discuss the general theory of rational maps.

Let $V$ and $W$ be varieties over $k$, and consider pairs $(U, \varphi_U)$
where $U$ is a dense open subset of $V$ and $\varphi_U$ is a regular map
$U \to W$. Two such pairs $(U, \varphi_U)$ and $(U', \varphi_{U'})$ are said
to be \textbf{equivalent} if $\varphi_U$ and $\varphi_{U'}$ agree on
$U \cap U'$. An equivalence class of pairs is called a \textbf{rational map}
$\varphi \colon V \dashrightarrow W$. A rational map $\varphi$ is said to be
\textbf{defined} at a point $v$ of $V$ if $v \in U$ for some
$(U, \varphi_U) \in \varphi$. The set $U_1$ of $v$ at which $\varphi$ is
defined is open, and there is a regular map $\varphi_1 \colon U_1 \to W$ such
that $(U_1, \varphi_1) \in \varphi$ --- clearly,
$U_1 = \bigcup_{(U, \varphi_U) \in \varphi} U$ and we can define $\varphi_1$
to be the regular map such that $\varphi_1 | U = \varphi_U$ for all
$(U, \varphi_U) \in \varphi$.

The following examples illustrate the major reasons why a rational map
$V \dashrightarrow W$ may fail to extend to a regular map on the whole of
$V$.

\begin{itemize}
\item[(a)] Let $W$ be a proper open subset of $V$; then the rational map
  $V \dashrightarrow W$ represented by $\mathrm{id} \colon W \to W$ will not
  extend to $V$. To obviate this problem, we should take $W$ to be complete.
\item[(b)] Let $C$ be the cuspidal plane cubic curve $Y^2 = X^3$. The
  regular map $\mathbb{A}^1 \to C$, $t \mapsto (t^2, t^3)$, defines an
  isomorphism $\mathbb{A}^1 \smallsetminus \{0\} \to C \smallsetminus \{0\}$.
  The inverse of this isomorphism represents a rational map
  $C \dashrightarrow \mathbb{A}^1$ which does not extend to a regular map on
  $C$ because the map on function fields doesn't send the local ring at
  $0 \in \mathbb{A}^1$ into the local ring at $0 \in C$. Roughly speaking, a
  regular map can only map a singularity to a worse singularity. To obviate
  this problem, we should take $V$ to be nonsingular (in fact, nonsingular
  is no more helpful than normal).
\item[(c)] Let $P$ be a point on a nonsingular surface $V$. It is possible
  to ``blow-up'' $P$ and obtain a surface $W$ and a morphism
  $\alpha \colon W \to V$ which restricts to an isomorphism
  $W \smallsetminus \alpha^{-1}(P) \to V \smallsetminus P$ but for which
  $\alpha^{-1}(P)$ is the projective line of ``directions'' through $P$. The
  inverse of the restriction of $\alpha$ to
  $W \smallsetminus \alpha^{-1}(P)$ represents a rational map
  $V \dashrightarrow W$ which does not extend to all $V$, even when $V$ and
  $W$ are complete --- roughly speaking, there is no preferred direction
  through $P$, and hence no obvious choice for the image of $P$.
\end{itemize}

In view of these examples, the next theorem is best possible.

\paragraph{Theorem 3.1.}
\textit{A rational map $\varphi \colon V \dashrightarrow W$ from a normal
variety $V$ to a complete variety $W$ is defined on an open subset $U$ of
$V$ whose complement $V \smallsetminus U$ has codimension $\ge 2$.}

\paragraph{Proof.}
Assume first that $V$ is a curve. Thus we are given a nonsingular curve $V$
and a regular map $\varphi \colon U \to W$ from an open subset of $V$ which
we want to extend to $V$. Consider the maps
\[
  \begin{array}{ccc}
      &            & V \\
      & \nearrow   & \Big\uparrow \text{project} \\
    U & \;\lhook\joinrel\xrightarrow{\;u \mapsto (u, \varphi(u))\;}\;
        & Z \subset V \times W \\
      & \underset{\varphi}{\searrow}
        & \Big\downarrow \text{project} \\
      &            & W.
  \end{array}
\]
% [diagram: hooked horizontal arrow U ↪ Z ⊂ V × W labelled "u ↦ (u, φ(u))";
%  unlabelled diagonal arrow U → V; diagonal arrow U → W labelled φ;
%  vertical arrows Z → V (upward) and Z → W (downward), each labelled "project"]
Let $U'$ be the image of $U$ in $V \times W$, and let $Z$ be its closure.
The image of $Z$ in $V$ is closed (because $W$ is complete), and contains
$U$ (the composite $U \to V$ is the given inclusion), and so $Z$ maps onto
$V$. The maps $U \to U' \to U$ are isomorphisms. Therefore,
% [proof of Theorem 3.1 continues on the next page]
